\section{Conclusão}
\label{sec:conclusao}

Mineramos as quatro métricas DORA em 300 repositórios open-source populares que usam GitHub
Actions, ao longo da janela anual de 2025. A amostra revelou desenvolvimento organizado em
ciclos discretos --- mediana de 0,46 release por semana e lead time por release de 324,9 h
(13,5 dias) --- desafiando a premissa de entrega contínua diária herdada do contexto
corporativo; a maior parte concentra-se em \emph{Medium} (154) e \emph{High} (110), com apenas
9 em \emph{Elite}. A estabilidade percebida varia drasticamente conforme a lente: o
\emph{proxy} de CI sugere maturidade \emph{Elite} (5,63\%), o de entrega indica correção
frequente (26,92\%), e os dois discordam de faixa em 72,3\% dos repositórios.

O achado central não é nenhum desses valores, e sim a sua dependência de escolhas que o
método não torna visíveis. Até 31,3\% dos repositórios migram de categoria entre combinações
e --- graças ao desenho que separa os dois efeitos --- trocar a \textbf{variante da métrica}
desloca a classificação quase quatro vezes mais do que trocar a \textbf{unidade de deploy}
(28,3\% contra 7,3\%). Somado a isso, 86,3\% mudam de categoria se uma única métrica avançar um
degrau, porque as faixas discretizam medidas contínuas sem sinalizar quem está na fronteira.

Decorrem daí três consequências práticas: relatar uma classificação DORA sem explicitar as
definições adotadas não é informativo; o CFR calculado sobre falhas de CI não é comparável às
faixas do relatório DORA, que tratam de produção; e comparações de categoria entre estudos ou
execuções precisam reportar quantos repositórios estavam próximos de um limite, sob pena de
atribuir a erro de cálculo o que é efeito da discretização --- lição que a replicação cruzada
nos impôs na prática.

Como trabalhos futuros, propõem-se validar a heurística corretiva em uma segunda amostra
rotulada, independente daquela em que foi calibrada; estender a observação para além de uma
janela anual, permitindo falar de tendência; e substituir as faixas discretas por uma medida
contínua de desempenho, que não reproduza a fragilidade de fronteira documentada aqui.

\vspace{0.2cm}
\noindent\textbf{Disponibilidade.} Código, dados, dicionário de dados e todas as tabelas e
figuras deste artigo estão em
\url{https://github.com/Paulooassis/Laboratorio-03-medicao-de-software}, reexecutáveis com
\texttt{python -m analise}. O acompanhamento de Issues e sprints está em
\url{https://github.com/Paulooassis/Laboratorio-03-medicao-de-software/projects}.
